\chapter{Introduction}

In this work, we will address some of the challenges associated with using Raman Spectroscopy (RS) for the task of glioma tumor grading \cite{Trevisi2025,stupak2025raman}. Our goal is to evaluate the usability of RS in a clinical setting, and for this we consider Machine Learning (ML) models for classification of glioma tumors based on RS spectra, and we pay considerable attention to the statistical validation framework, implementing a robust cross-validation scheme to ensure that the results are scientifically valid and reproducible; we consider augmentation as a way to address the small sample size, and to accomplish this we develop a generative model to generate synthetic RS spectra that are similar to real spectra. We then test the ML models in various regimes and interpret the results.

RS is a molecular fingerprinting technique that can be used to identify the chemical composition of a sample \cite{Jones2019}. It has many applications in materials science, chemistry, and biology, and has been shown to be a promising tool for medical diagnostics. In particular, RS has been used to identify cancerous tissue in various organs, including the brain \cite{Auner2018}.
The current medical standard for glioma grading is histopathological testing, which involves laboratory processing of the sample, a technique that requires days or even weeks, therefore it is not suitable for intraoperative diagnosis \cite{whitfield2022classificationgliomas}.
In contrast, RS can be applied on the raw biopsy sample, offering near instant diagnostics \cite{stupak2025raman, Trevisi2025}. If RS can be shown to be a reliable method for glioma grading, it could be used to provide real-time feedback to surgeons during operations, potentially improving the speed of medical procedures and patient outcomes.

Despite the potential of RS for glioma grading, there are several important limitations that make the adoption of RS in clinical practice challenging. The first limitation is the lack of large, publicly available datasets of RS spectra from glioma tumors. This makes it difficult to train machine learning models for glioma grading, as these models require large amounts of data to learn from. The second limitation is the high variability in RS spectra between different patients and even between different samples from the same patient. There is a strong correlation of spectra collected in a spectroscopy session, but high variability exists between sessions. This variability can make it difficult to develop a model that generalizes well to new data. The third limitation is the high dimensionality of RS spectra, and low number of samples. RS spectra are typically collected at hundreds or thousands of different wavelengths, resulting in high-dimensional data. However, the number of available samples is often limited to a few dozen or hundred, which can make it difficult to train machine learning models that generalize well to new data \cite{wang2024current, shen2026advances}.

Because of these limitations, the technique is still in its research phase, and is not yet widely used in clinical practice. Moreover, the variability of RS spectra acquired from different research groups and different instruments makes it difficult to compare results across studies.

In this work, we aim to address some of these limitations by using synthetic data to augment the available datasets of RS spectra from glioma tumors. We propose a Conditional Variational Autoencoder (CVAE) architecture to generate synthetic RS spectra that are similar to real spectra, and we will use these synthetic spectra to train machine learning models for glioma grading. We will then evaluate the performance of these models on real RS spectra, and compare their performance to models trained only on real data.

\section*{Contributions}



Our contributions to the research of the applicability of RS in the context of glioma grading are the following

\begin{itemize}
  \item \emph{A spatio-spectral annotation pipeline} for separating tumour spectra from
        background and non-tissue regions on a per-slide basis, combining median filtering,
        NMF, and GMM clustering under visual inspection
        (Section~\ref{sec:preprocessing_steps}). Different decomposition and clustering methods exists, and other methods have been used for this task, for example PCA and DBSCAN was used as a preprocessing step by Lita et al. \cite{apollonov2024raman}. However some researchers argue that two processes involved in acquiring a Raman spectrum: light scattering and sample molecular mix are both non-negative and additive in nature. Light scattering involves measuring light intensity, where all intensities are greater than or equal to 0, and the chemical composition mixing suggests that at least partly or at least in some situations, the measured spectrum could be a linear combination of non-negative component spectra corresponding to the constituenting molecules of the sample \cite{Alix2022, Jafari2025}. Hence, we make the case that NMF might allign better for this task. The visual inspection step could be a benefit as one could better control the quality of spectra selected for processing, but it could also induce unintended selection bias.

  \item \emph{A conditional generative model for Raman spectra}, a 1D-convolutional
        $\beta$-CVAE trained under a multi-task objective combining reconstruction, KL
        regularisation, and auxiliary classification (Section~\ref{sec:cvae_architecture}).

  \item \emph{An evaluation under a strict patient-level protocol}, comparing TR/TR, TS/TR
        and TSR/TR regimes across four downstream classifiers and two tasks
        (Chapter~\ref{sec:experimental_design}), with uncertainty estimated by a slide-level
        cluster bootstrap rather than at the spectrum level
        (Section~\ref{sec:statistical_eval}). The two protocols deliberately use
        different units: splitting at the patient level, which is the stricter
        requirement for preventing leakage, and resampling at the level of the
        acquisition, which is the cluster that governs spectral variability. This methodology is required for faithful estimation of real world performance of the models. We treat each sample as the unit of test and not individual spectra.

  \item \emph{Classification by Reconstruction}, an inference strategy that classifies a
        spectrum by which class-conditioned reconstruction best matches it, using no
        discriminative classifier at all (Section~\ref{sec:cbr}).

  \item \emph{A characterisation of the domain gap}, showing that the generator reproduces
        class-mean spectra faithfully but collapses their variance, and that this fidelity /
        diversity split explains the observed TS/TR and TSR/TR behaviour
        (Chapter~\ref{sec:discussion}).
\end{itemize}

\section*{Outline}

The rest of this thesis is divided as follows:

Chapter 2 reviews Raman Spectroscopy, its use in medical diagnostics, glioma grading, and synthetic data
for spectroscopy. We cover the fundamental physics process of RS and highlight its use cases, particularly in medicine, and more specifically in tumour grading of gliomas.

Chapter 3 introduces the theoretical mathematical background of the tools used in this work. It describes each tool and mentions how it is applied in the particular context of this work.

Chapter 4 describes the methodology as a pipeline of the tools used and processing steps applied to the data. This chapter acts as a blueprint of the exact tools and parameters used from preprocessing to the generator definition and training.

Chapter 5 defines the evaluation protocols and statistical framework. It describes the design of the statistical tests, using a clustering sampling and bootstrap design required for this type of data.

Chapter 6 summarizes the main results, while pointing to the extended tables of results in the appendinx.

Chapter 7 interprets the results and states limitations and addresses possible future work.

Chapter 8 concludes the work, highlighting the main insights and limitations of this study.

